% ---------------------------------------------------------------
\section{Limitations and Future Work}\label{sec:limitations}
% ---------------------------------------------------------------

The Bailout Protocol's guarantees---SAFETY, LIVENESS, and ACCOUNTABILITY---are derived from structural invariants enforced by the \fact{} layer's pre-commit logic. These invariants hold unconditionally within the system's stated assumptions: a computationally intact parent-pointer chain $\alpha(n)$, a non-adversarial communication substrate, and a reachable human anchor $h(n)$. This section identifies the environmental constraints that bound the protocol's operational scope, classifies each as fundamental or addressable, and charts the research agenda motivated by each.

% --- 6.1 Limitations ---
\subsection{Limitations}\label{sec:limitations:items}

\subsubsection{Human Response Latency}
\label{sec:lim:latency}

The LIVENESS guarantee is bounded above by the human-response timeout $T_{\text{human}}$, defined in the state-machine specification (Section~\ref{sec:bailout-protocol}). When $h(n)$ fails to respond within $T_{\text{human}}$, the state machine transitions to \texttt{TIMEOUT} and the originating node enters a quiescent error state. SAFETY is preserved---consequential action stalls---but the exception remains unresolved until $h(n)$ regains availability or the deployment enters manual remediation.

The constraint is compounded in globally distributed deployments where $h(n)$ occupies a different timezone, regulatory jurisdiction, or operational context. The protocol currently provides no mechanism to distinguish between an \textit{unavailable} anchor (temporary loss of platform access, medical emergency) and an \textit{unresponsive} anchor (cognitive overload, disengagement, or adversarial compromise of authentication credentials). These represent categorically different failure modes requiring materially different responses: pool reassignment or security intervention in the latter case, simple availability wait in the former. Conflating them delays incident response and may extend quiescent error periods beyond what the operational context can tolerate.

Time-sensitive domains---clinical decision support, real-time infrastructure control, autonomous transport, financial settlement---cannot guarantee latency-bounded escalation without modifications to the anchor designation model. In these domains, $T_{\text{human}}$ is often smaller than the minimum time required for a human to formulate a calibrated directive under novel exception conditions, creating a structural tension between timeliness and response quality that is not resolvable by tuning $T_{\text{human}}$ alone. We classify this limitation as \emph{fundamental under the single-anchor model} and \emph{addressable under the distributed anchor pool model} (Section~\ref{sec:fw:anchor-pools}).

\subsubsection{Adversarial Conditions}
\label{sec:lim:adversarial}

The Compulsory Human Escalation theorem (Section~\ref{sec:guarantees-theorem}) assumes a computationally intact parent-pointer chain $\alpha(n)$ and a non-adversarial communication substrate. Both assumptions are violated under active adversarial conditions. A compromised parent node $p$ in the propagation chain may exhibit Byzantine behavior: it may acknowledge escalation messages without forwarding them, selectively drop \texttt{BAIL\_PENDING} notifications, or inject fabricated acknowledgments that satisfy the \fact{} layer's entry format while masking the absence of a genuine human response. The parent pointer $\alpha(n)$ remains structurally immutable (enforced by the \fact{} layer), but the node at that pointer position can behave arbitrarily.

The Pathway Health Registry (Section~\ref{sec:bailout-state-machine}) partially mitigates degraded parent nodes by excluding them from pathway selection after a configurable health-check threshold. However, the health check is a network operation that a sufficiently capable adversary controlling the communication substrate can interfere with or fabricate. The adversary does not need to compromise any individual node; controlling the network transport layer suffices to make the health registry report false positives and mask the absence of genuine message propagation.

The protocol provides no cryptographic attestation of propagation-path integrity: it cannot guarantee that \texttt{escalation-sent} Forensic Ledger entries correspond to genuinely received and forwarded messages rather than locally fabricated acknowledgments. This gap is fundamental when the adversary controls the network layer rather than individual nodes. Closing it requires a chain-of-cryptographic-receipts extension with per-hop latency and storage overhead proportional to path length.

The diagnostic context attached to each escalation---exception classification, Bounds Engine reasoning trace, and resolution history---additionally exposes an information surface that a compromised propagation-chain node can exploit. Exception classifications reveal the \bounds{} layer's detection logic; resolution histories expose which failure patterns have been encountered; the Bounds Engine's reasoning trace may expose sensitive business logic, spawn-tree structure, or the principal's operational patterns. A systematic probing adversary can use diagnostic responses to build a detailed internal architecture map and plan targeted exploitation accordingly. We classify the adversarial-condition gap as \emph{fundamental under the single-path propagation model} and \emph{addressable under a multi-path propagation model with cryptographic attestation} (Section~\ref{sec:fw:verification}).

\subsubsection{Scalability}
\label{sec:lim:scalability}

The protocol's accountability model is structurally single-threaded at the human anchor layer: each node $n$ designates a single anchor $h(n)$, and all unresolved exceptions route to that one principal. When $N$ simultaneous nodes each trigger the Bailout Protocol, $h(n)$ receives $N$ independent escalations with full diagnostic context. The protocol provides no mechanism for prioritization, batching, automatic triage, or selective compression of concurrent exceptions. All escalations compete equally for the anchor's attention.

The problem is not merely throughput. The single-anchor model creates a cognitive overload condition at high concurrency: even if $h(n)$ could theoretically process $N$ escalations, human cognitive bandwidth for exception evaluation is bounded, and the protocol provides no schema for comparative risk assessment across concurrent escalations. A low-severity capability-bound violation, a medium-severity timeout, and a high-severity norm violation carry identical escalation formats in the current specification. The anchor must perform triage without structural support from the protocol itself.

The forensic record additionally grows linearly with the number of exceptions processed across all nodes. The SHA-256 hash-chained Forensic Ledger imposes storage and verification overhead at scale: hash-chain verification time is $O(M)$ in the number of entries $M$, meaning that verification of a long-lived ledger becomes progressively more expensive. The current specification provides no tiering, summarization, selective pruning, or checkpoint-based short-circuiting strategy for routine low-severity exceptions. We classify scalability as \emph{addressable through architectural extensions}: distributed anchor pools address the triage bottleneck (Section~\ref{sec:fw:anchor-pools}), and tiered Forensic Ledger management is a natural extension of the existing hash-chain structure.

\subsubsection{Exploitation Risk}
\label{sec:lim:exploitation}

A motivated adversary who can induce exceptions at a target \bxthree{} node can mount a denial-of-service escalation attack: by repeatedly triggering condition TC, the adversary forces repeated \texttt{BOUNDED} states and eventual \texttt{BAIL\_PENDING} transitions. Each transition consumes the human anchor's evaluation bandwidth, generates Forensic Ledger entries, and may force the node into quiescent error state. The protocol provides no rate-limiting mechanism based on the source or pattern of exception triggers.

This risk is acute because TC depends in part on the Bounds Engine's self-assessed confidence threshold $\tau$. An adversary who crafts inputs that systematically reduce the Bounds Engine's confidence below $\tau$ can induce pathological escalation from a well-designed node. While the \fact{} layer prevents machine actors from suppressing escalations once generated, it does not constrain the generation of escalation conditions, and the Bounds Engine's trigger evaluation is treated as a black box without external auditability of TC firing rates. The escalation interface's exposure of diagnostic context additionally enables a probing adversary to refine their attack surface: each diagnostic response from a failed escalation reveals which Bounds Engine categories are active, which thresholds are configurable, and which resolution patterns the anchor has accepted in prior similar cases.

We classify exploitation risk as \emph{addressable through formal TC analysis and rate-limiting extensions} (Section~\ref{sec:fw:verification}).

% --- 6.2 Future Work ---
\subsection{Future Work}\label{sec:future-work}

The research agenda motivated by the limitations above is organized around three independent but practically interdependent axes: formal verification of the protocol's guarantees under adversarial conditions, adaptive adjustment of timing parameters to operational conditions, and distribution of the human anchor function to eliminate single points of failure.

\subsubsection{Formal Verification}
\label{sec:fw:verification}

The Bailout Protocol's invariants (INV-1 through INV-3) and the Compulsory Human Escalation theorem are currently stated as pen-and-paper propositions derived from the state-machine definition and \fact{} layer enforcement semantics. While the reasoning is internally consistent, pen-and-paper proofs are subject to subtle logical gaps that mechanized verification is designed to eliminate. A natural next step is mechanized formal verification using TLA+\ \cite{lamport2019tlaplus} or Coq \cite{coq2024}.

Target verification properties include: (i) the \textit{Compulsory Human Escalation theorem}, expressed as a temporal logic formula asserting that for all nodes $n$ and all exception states $e$, the protocol reaches $h(n)$ within a bounded number of transitions; (ii) INV-1, asserting that no node can be in an active-bailout state without the state machine being in \texttt{ESCALATING}; and (iii) INV-2, asserting the timeout-forced nature of the \texttt{BAIL\_PENDING} $\rightarrow$ \texttt{ESCALATING} transition. Mechanized verification of INV-2 is particularly important because it establishes that the transition to \texttt{ESCALATING} is structurally forced, not operationally contingent.

Mechanized verification would additionally enable rigorous analysis of the protocol under Byzantine fault conditions. A formal model could establish the conditions under which the Compulsory Human Escalation theorem continues to hold when up to $f$ parent nodes in the propagation chain exhibit arbitrary failure modes. This requires extending the spawn-tree model to include redundant propagation pathways and majority-voting semantics at the \fact{} layer---the same extension needed for distributed human anchor pools. A TLA+ specification of the \bxthree{} spawn tree and Bailout Protocol state machine is planned as a companion open-reference artifact for researchers and practitioners.

A second verification track addresses the propagation-path integrity gap identified in Section~\ref{sec:lim:adversarial}. Cryptographic attestation of escalation propagation---where each parent node on the path provides a chain-of-custody receipt attesting to genuine receipt and forwarding rather than local fabrication---can be formalized as a property analogous to delivery receipts in certified messaging systems. A formal model of this property would define the minimal cryptographic primitives required to make fabricated acknowledgments detectable, establish the tradeoff between attestation strength and per-hop latency overhead, and provide the specification basis for the multi-path propagation extension.

A third verification track models the TC trigger condition under adversarial input manipulation to establish the conditions under which the Bounds Engine's confidence threshold $\tau$ can be systematically degraded. This requires a threat model for the input surface of the Bounds Engine, a formal analysis of the trigger condition's sensitivity to input perturbation, and a specification of the rate-limiting mechanism required to bound the adversarial exploitation rate.

\subsubsection{Adaptive Timeout}
\label{sec:fw:adaptive-timeout}

The fixed timeout parameters $T_{\text{bounds}}$, $T_{\text{prop}}$, $T_{\text{cap}}$, and $T_{\text{human}}$ are provisioned at node initialization and held constant for the node's operational lifetime. This design simplifies implementation and makes timing behavior statically predictable, which is advantageous for certification and audit contexts. However, static timeouts cannot respond to operational realities: network latency varies with time of day and geographic dispersion; compute load variations affect processing time inside node reasoning cycles; anchor availability varies with timezone and work schedule.

We intend to investigate an adaptive timeout extension in which each parameter is replaced by a feedback-controlled function that adjusts based on observed historical values. Specifically, let $T_{\text{prop}}^{\text{obs}}(n)$ denote the empirically observed propagation latency from node $n$ to $h(n)$ across the last $k$ successful escalations. A weighted moving average $\hat{T}_{\text{prop}}(n) = \alpha \cdot T_{\text{prop}}^{\text{obs}}(n) + (1-\alpha) \cdot \hat{T}_{\text{prop}}(n-1)$ replaces the static $T_{\text{prop}}$ provisioning, with $\alpha \in (0, 1)$ controlling adaptation speed. The \fact{} layer records both the provisioned and observed values, maintaining verifiability: an adaptive timeout does not compromise the forensic record because the \fact{} layer continues to log the actual observed values alongside the provisioned parameters, and the timeout transition triggers remain deterministic functions of observed rather than estimated values.

The benefit of adaptive timeout is improved liveness in non-adversarial deployments: a node operating in a low-latency environment escalates faster without waiting for a conservatively provisioned fixed timeout, while a node in a high-latency environment automatically extends its propagation window to accommodate realistic round-trip times.

\subsubsection{Distributed Human Anchor Pools}
\label{sec:fw:anchor-pools}

The single-anchor designation $h(n)$ is the root cause of three of the four limitations described above: the latency ceiling (unavailable single anchor), the scalability bottleneck (single-threaded triage), and the exploitation vulnerability (focused denial-of-service bandwidth). This motivates the investigation of a distributed human anchor pool architecture as the primary structural extension to the protocol.

Under the pool model, rather than designating a single principal $h(n)$ for each node $n$, the protocol designates a pool $H(n) = \{h_1, h_2, \ldots, h_m\}$ of $m$ human anchors, each with authenticated access to the \fact{} layer's escalation endpoint. The protocol requires acknowledgment from any $k$-subset of $H(n)$ to transition from \texttt{ESCALATING} to \texttt{HUMAN\_ACKNOWLEDGED}, providing Byzantine fault tolerance against up to $m - k$ unresponsive or adversarial anchors. The values of $m$ and $k$ are provisioned at node initialization and recorded in the \fact{} layer; they are immutable once set for a given node (changing the pool requires a new \fact{} layer entry, not an overwrite).

The pool architecture addresses all three single-anchor failure modes. Availability improves because the system tolerates the absence of any individual anchor without blocking escalation. Scalability improves because the triage burden is distributed: a high-severity exception can be routed to the first available anchor rather than waiting for a designated principal to become reachable. The denial-of-service attack vector is widened but its per-anchor cost is reduced: an adversary must exhaust the attention of $k$ independent anchors to achieve the same effect that exhausting a single anchor achieved under the original model.

The pool architecture introduces open design questions requiring systematic resolution. First, directive-issuance semantics must handle the case where different anchors in the pool issue conflicting directives for the same escalation. A majority-voting scheme could resolve this but requires the \fact{} layer to support multi-author directive entries and introduces non-determinism in the resolution path that must be accounted for in the forensic record. Second, dynamic pool reconfiguration must not invalidate in-flight escalations: an anchor that departs mid-escalation (voluntary exit, credential revocation, adversarial compromise) must not cause the protocol to stall or deadlock. A minimum-active-anchors threshold below which no new escalations can be initiated, with in-flight escalations continuing with the remaining anchors, is one structural approach. Third, pool membership at the time of each escalation must be recorded in the Forensic Ledger, and retroactive changes to pool membership must not affect the verifiability of past escalation records. Fourth, pool sizing and $k$-threshold calibration for different deployment contexts requires empirical validation that the current single-anchor deployment cannot support.

% ---------------------------------------------------------------
\section{Conclusion}\label{sec:conclusion}
% ---------------------------------------------------------------

The Bailout Protocol addresses a foundational problem in accountable autonomous systems: the absence of an architecturally enforced mechanism that compels unresolved exception states to reach a designated human principal, bypassing all machine actors in the intervening hierarchy. We have shown that this compulsion is not achieved through policy or convention---the two approaches that characterize the majority of HITL systems in current deployment---but through the structural integration of the protocol's state machine into the \fact{} layer's pre-commit enforcement substrate. The integration makes it impossible for any machine actor, including the agent whose behavior triggered the exception, to suppress, redirect, or substitute the human escalation path. This is the protocol's core claim, and it is enforced architecturally rather than argued operationally.

The protocol makes three concrete contributions. First, it provides a precise operational semantics for exception escalation in bounded-autonomy agentic systems, expressed as a deterministic state machine $\mathcal{B}$ with formally specified invariants (INV-1 through INV-3) that encode the protocol's core safety guarantees and a clearly defined trigger condition TC that governs activation. The state-machine specification is self-contained: its properties are derived from its own structure and from the \fact{} layer's enforcement semantics, not from assumptions about agent behavior. Second, it establishes the \textit{Compulsory Human Escalation theorem}: for any node $n$ in a \bxthree{} spawn tree and any exception state $e$ at $n$, the protocol's propagation path $C(n)$ reaches $h(n)$ in finite time, independent of the reasoning strategy employed by any machine actor in the chain. The theorem's guarantee is structural, not probabilistic; it follows from the immutability of the parent pointer and the finiteness of the spawn tree depth, both enforced by the \fact{} layer. Third, it demonstrates that the upstream accountability guarantee of the \bxthree{} Framework---that systems fail \emph{upward} into human consciousness, never \emph{downward} into autonomous chaos---is not a design aspiration but an architecturally enforced operational property, made compulsory by the three-layer model and the pre-commit enforcement mechanism.

The \bxthree{} Framework's role is constitutive of all three contributions, not merely incidental. The three-layer model provides the structural substrate on which the protocol operates: the \purpose{} layer defines the human anchor $h(n)$ and the node's authorized operational scope; the \bounds{} layer houses the Bounds Engine that detects and classifies constraint violations and evaluates TC; and the \fact{} layer enforces the state machine transitions that drive the protocol, records every escalation event in the immutable Forensic Ledger, and prevents any machine actor from intercepting or modifying the escalation path. Without this three-layer separation, the Bailout Protocol would reduce to a discretionary interface whose enforcement depended on the correct behavior of the very machine actors it is designed to govern. The \bxthree{} Framework is not a container for the Bailout Protocol; it is the mechanism by which the protocol becomes compulsory.

We acknowledge that the protocol operates within real constraints: human response latency imposes a practical upper bound on escalation timeliness that is fundamental under the single-anchor model; adversarial network conditions can compromise the propagation chain in ways the current single-path specification cannot detect or resist; single-anchor designation limits scalability and creates a focused exploitation surface; and the escalation interface exposes diagnostic context that a compromised node can observe. These are not reasons to abandon the approach---they are the precise targets of the formal verification, adaptive timeout, and distributed human anchor pool research programs outlined in Section~\ref{sec:future-work}. Each limitation corresponds to an open research question with a clear resolution path, and each resolution path is architecturally compatible with the protocol's existing invariants.

The broader significance of the Bailout Protocol extends beyond the \bxthree{} Framework's specific architecture. The protocol demonstrates that the upstream accountability guarantee---a long-standing design principle in human-in-the-loop AI literature, traceable to Wiener's original conception of the human as the final corrective authority \cite{wiener1954human}---can be rendered formally compulsory through architectural enforcement rather than left to operational convention. The structural principles underlying the protocol---state-machine embedding in a tamper-evident enforcement layer, immutable human-anchor designation, and deterministic timeout-forced escalation---are applicable to accountable autonomous system design more broadly, independent of the specific reasoning architecture or domain of deployment. Any autonomous system that issues consequential actions without requiring human review of the decision process is, by design, incapable of failing safely. The Bailout Protocol shows what failing safely actually requires: not a human in the loop, but a structural compulsion that makes the human the unavoidable terminus of every unresolved exception.
